\subsection{Exact ordinary lifts for cyclic and two-group outer images}

\begin{statement}{D.34}{Proved}{Cyclic outer image admits every multiplier-one peripheral base automorphism}
Let
\[
 1\longrightarrow A\longrightarrow E
       \xrightarrow{\pi}D\longrightarrow1
 \tag{D.34.1}
\]
be a finite extension with $Z(A)=1$ and
a given source epimorphism
$q_E:\widehat F\twoheadrightarrow E$.
Let $X_i=q_E(x_i)$ and $d_i=\pi(X_i)$.
Write
$\epsilon:D\to\operatorname{Out}(A)$
for its actual outer action, and suppose
its image is cyclic.
Every $\alpha\in\operatorname{Aut}(D)$
with $\alpha(d_i)\in d_i^D$ has an
actual automorphism lift
$\widetilde\alpha\in\operatorname{Aut}(E)$
satisfying
$\widetilde\alpha(X_i)\in X_i^E$.
The full group $\operatorname{Out}(A)$
need not be abelian.
\end{statement}

\begin{proof}
As in \textup{(D.33.3)}, use the actual
centerless pullback isomorphism
\[
 E\xrightarrow{\sim}
       \operatorname{Aut}(A)\times_{\operatorname{Out}(A),\epsilon}D,
 \qquad e\longmapsto(\operatorname{Ad}_e,\pi(e)).
 \tag{D.34.2}
\]
Choose $h_i\in D$ with
$\alpha(d_i)=h_i d_i h_i^{-1}$.
All $\epsilon(h_i)$ and
$\epsilon(d_i)$ lie in the cyclic
image $C$ and commute.
Since $d_1,d_2$ generate $D$,
this proves
\[
                         \epsilon\alpha=\epsilon.
 \tag{D.34.3}
\]
Choose a generator $t$ of $C$,
of order $n$, and write
$\epsilon(d_i)=t^{a_i}$ and
$\epsilon(h_i)=t^{b_i}$.
The actual source relation gives
$a_1+a_2+a_3=0\pmod n$, and
$a_1,a_2$ generate $\mathbf Z/n$.

There is an $r\in\mathbf Z/n$ with
\[
                 r-b_i\in a_i(\mathbf Z/n)
                         \qquad(i=1,2,3).
 \tag{D.34.4}
\]
Here is a constructive proof.
For every $p^e$ exactly dividing
$n$, at least one of $a_1,a_2$
is a unit modulo $p$, and the
sum relation makes at least two
of $a_1,a_2,a_3$ units modulo $p$.
Thus at most one index imposes
a congruence on $r$ in that
primary component.
For that index, take
$r=b_i\pmod {p^e}$; if none
does, take $r=0\pmod {p^e}$.
The Chinese remainder isomorphism
combines these choices and proves
\textup{(D.34.4)}.
Solving the three resulting
linear congruences gives integers
$k_i$ with
$a_i k_i=r-b_i\pmod n$.

Put $h_i'=h_i d_i^{k_i}$.
Right multiplication by a power
of $d_i$ retains its conjugacy
equation, and
\[
                  \epsilon(h_i')=t^r
                           \qquad(i=1,2,3).
 \tag{D.34.5}
\]
Choose an actual
$c\in\operatorname{Aut}(A)$
representing the outer class
$t^r$.  In \textup{(D.34.2)}
define
\[
       \widetilde\alpha(g,d)
                    =(cgc^{-1},\alpha(d)).
 \tag{D.34.6}
\]
Conjugation by the outer class
of $c$ fixes every element of
$\epsilon(D)=C$, and
\textup{(D.34.3)} holds.
Thus \textup{(D.34.6)} preserves
the actual pullback condition.
Its inverse is
$(g,d)\mapsto(c^{-1}gc,\alpha^{-1}(d))$,
so it is an automorphism over
$\alpha$.
If $X_i=(g_i,d_i)$, then
$(c,h_i')$ belongs to the
pullback by
\textup{(D.34.5)}, and
\[
       (c,h_i')X_i(c,h_i')^{-1}
           =(cg_ic^{-1},\alpha(d_i))
           =\widetilde\alpha(X_i).
 \tag{D.34.7}
\]
This is the required ordinary
peripheral lift, transported
through \textup{(D.34.2)}.
\end{proof}

\begin{statement}{D.35}{Proved}{A two-group outer image admits every order-five multiplier-one peripheral shadow}
In \textup{(D.34.1)}, retain
$Z(A)=1$ and the given
two-generated source.
Suppose the actual image
$\epsilon(D)$ is a two-group.
Every order-five automorphism
$\alpha$ of $D$ preserving
the three labelled peripheral
conjugacy classes has an
ordinary peripheral automorphism
lift of order five to $E$,
fixing $A$ pointwise.
In particular an order-five
ordinary obstruction with
kernel $A_6^\Omega$ requires
its transitive factor-permutation
image to have non-two-power
order.  It cannot occur with
two simple factors or with
any purely two-group
factor-permutation action.
\end{statement}

\begin{proof}
Let $P$ be the characteristic
maximal two-group quotient
of the finite group $D$.
It is the quotient by the
intersection of kernels of
all maps to finite two-groups.
There are only finitely many
distinct such kernels in $D$;
the resulting quotient is a
subgroup of their finite
product and is itself a
two-group.  The given source
shows that $P$ is generated
by at most two elements.

There is no nontrivial
automorphism of order five
of a two-generated finite
two-group $P$.
To prove this directly,
let $\Phi(P)$ be its
Frattini subgroup, and put
$d=\dim_{\mathbf F_2}P/\Phi(P)\leq2$.
The action on this quotient
has image in
$\operatorname{GL}_d(\mathbf F_2)$,
whose order divides six.
The kernel of this action
is a two-group.
Indeed fix one ordered basis
of $P/\Phi(P)$ and consider
all ordered generating tuples
lifting it.  Every lift is
generating: a proper generated
subgroup would lie in a
maximal subgroup of index two,
contradicting its full image
modulo $\Phi(P)$.
There are exactly
$|\Phi(P)|^d$ such tuples.
An automorphism fixing the
Frattini quotient acts freely
on this set, because an
automorphism fixing a
generating tuple is the
identity.  Its group order
therefore divides the
two-power $|\Phi(P)|^d$.
This proves the claim,
including the trivial
quotient cases.

The order-five automorphism
$\alpha$ induces the identity
on $P$.  Since the actual
map $\epsilon$ factors
through $P$, it follows,
before any conjugator
adjustment, that
\[
                         \epsilon\alpha=\epsilon.
 \tag{D.35.1}
\]
Choose $h_i$ with
$\alpha(d_i)=h_i d_i h_i^{-1}$
and form
\[
              c_i=\alpha^4(h_i)\alpha^3(h_i)
                          \alpha^2(h_i)\alpha(h_i)h_i.
 \tag{D.35.2}
\]
Then $c_i\in C_D(d_i)$, and
\textup{(D.35.1)} gives the
literal power identity
\[
                         \epsilon(c_i)=\epsilon(h_i)^5.
 \tag{D.35.3}
\]
No commutativity of the
two-group image is assumed:
the five factors on the
right are the same element.
Choose an integer $k$ with
$5k\equiv1$ modulo the
two-power exponent of
$\epsilon(D)$ and put
$h_i'=h_i c_i^{-k}$.
It still conjugates $d_i$
to $\alpha(d_i)$ and now
satisfies $\epsilon(h_i')=1$.

In the actual centerless
pullback \textup{(D.34.2)},
the formula
\[
                     (g,d)\longmapsto(g,\alpha(d))
 \tag{D.35.4}
\]
is an automorphism by
\textup{(D.35.1)}.
It has order five, fixes
$A$ pointwise, and its
peripheral conjugators are
the actual elements
$(1,h_i')$ of the pullback.
This proves the lift.

For $A=A_6^\Omega$, the
exact outer group is
\textup{(D.33.9)}.
If its factor-permutation
image is a two-group,
then $\epsilon(D)$ is an
extension of a subgroup
of $(\mathbf F_2^2)^\Omega$
by a two-group and is
itself a two-group.
The result therefore
applies.  For two simple
factors, the transitive
permutation image is
the group of order two,
so that case is included.
\end{proof}

Both results retain the ordinary
centralizer freedom by explicit
maps.  Neither asserts the
corresponding lift with every
initially prescribed base
conjugator unchanged.
